\documentclass[10pt,a4paper]{amsart}
\usepackage[margin=19mm]{geometry}
\usepackage{amsmath,amssymb,mathtools}
\usepackage[scr=rsfs,cal=euler]{mathalfa}
\usepackage{xfrac}
\usepackage[unicode,hidelinks,bookmarks=false]{hyperref}
\newcommand{\ie}{i.e.\ }
\newcommand{\cf}{cf.\ }
\newcommand{\resp}{respectively\ }
\newcommand{\Subsection}{\subsection}
\providecommand{\nobreakdash}{\nobreak\hbox{-}}
\setlength{\emergencystretch}{2em}
\multlinegap0pt
\usepackage{bm}
\usepackage{amssymb}
\usepackage{cancel}
\usepackage[shortlabels]{enumitem}
\usepackage{xcolor}
\usepackage{enumerate}
\usepackage{float}
\usepackage{mathtools}
\newtheorem{theorem}{Theorem}
\newcommand{\bs}{\boldsymbol}
\title{St\"ackel and Eisenhart lifts, Haantjes geometry and Gravitation}
\author{Ond\v{r}ej Kub\r{u}, Piergiulio Tempesta}
\date{Source: arXiv:2509.19950v4 (CC BY 4.0)}
\begin{document}
\maketitle

\noindent\emph{Source and licence.} St\"ackel and Eisenhart lifts, Haantjes geometry and Gravitation. Version arXiv:2509.19950v4 (CC BY 4.0), distributed by arXiv under the Creative Commons Attribution 4.0 International licence, http://creativecommons.org/licenses/by/4.0/, as recorded in the arXiv licence metadata for this exact version and shown on the abstract page https://arxiv.org/abs/2509.19950v4. Full licence notice: \texttt{licenses/arxiv-2509.19950-CC-BY-4.0.txt}.\par\smallskip

\noindent\textit{Editorial scope.} Counted unit: the article's theorem theo:1 (source0.tex lines 609--638). The article's complete original proof of it is delivered with it (source0.tex lines 640--706, one \texttt{proof} environment of 67 lines). No mathematics is written, repaired or re-derived: the statement and the proof are the article's own lines, in the article's own notation, and no retained line is substituted or re-flowed.  Retained uncounted as the source-local context the proof needs: lines 445--447; lines 477--479.  Nothing was omitted from any retained span, so the package carries no omission record.  The retained lines cite 1 works, whose entries are transcribed verbatim from the article's own bibliography source in the e-print and delivered with short editorial labels recorded in the item's reference mapping.  No retained line contains a figure environment, an image-inclusion command or drawing code, so no image is delivered and the package declares no asset dependency.  Editorial formatting only, in addition: an AMS \texttt{amsart} preamble that restates the article's own declarations and macros used by the retained lines, namely the environments \texttt{theorem} on separate counters, together with the standard packages those lines need.  The retained environments print in this document as Theorem 1 in order of appearance, whereas the article prints the same items with the numbering of its own full document.  The complete native source of the article as distributed in the arXiv e-print for this exact version is bundled unchanged as \texttt{sources/185645/source0.tex}.
	\begin{equation} \label{eq:3.3}
		ds^2=\sum_{i,j=1}^{n} g_{ij}(\bs{q},t)dq^{i} d q^{j}-2V(\bs{q},t)dt^2+2\sum_{i=1}^{n} A_i(\bs{q},t)dq^i dt+2 dt du.
	\end{equation}

	\begin{equation}\label{eq:stackel_data}
		g_{j}^{kk}(\boldsymbol{q})=[(S^{(n)})^{-1}]_{j}^{k}(\bs{q}) \, , \qquad\qquad V_{j}(\bs{q})=\sum_{i=1}^n g^{ii}_j(\boldsymbol{q})\, W_i(q^i) \ .
	\end{equation}

\begin{theorem}[Stäckel lift and Riemannian Eisenhart lift]\label{theo:1}
	Let
	\begin{equation}\label{eq:H}
		H_j = \frac{1}{2}\sum_{i=1}^n g^{ii}_j(\boldsymbol{q})\, p_i^2 + V_j(\boldsymbol{q}), \qquad j=1,\ldots,n,
	\end{equation}
	be a Stäckel-separable completely integrable Hamiltonian system on $T^{*}Q$, separable in a coordinate chart
	$(q^1,\ldots,q^n)$ with associated Stäckel matrix $S^{(n)}$.
	Let the contravariant metric and Stäckel potentials be given by eqs.
	\eqref{eq:stackel_data}.
	Consider the lifting matrix
	\begin{equation}\label{eq:lifting-matrix}
		\mathcal{L}^{(n+1)}(\boldsymbol{q})=
		\begin{pmatrix}
			S^{(n)}(\boldsymbol{q}) & -\boldsymbol{W} \\
			\boldsymbol{0}^T & 1
		\end{pmatrix},
		\qquad
		\boldsymbol{W}=(W_1(q^1),\ldots,W_n(q^n))^T .
	\end{equation}
	Then, choosing the Stäckel functions $f_i = \frac{1}{2}p_i^2$ (for $i=1,\ldots,n$) and $f_{n+1} = p_u^2$, the Stäckel lift generated by $\mathcal{L}^{(n+1)}$ yields the independent lifted Hamiltonians on $T^{*}\tilde{Q}$, making the system completely integrable: 
	\begin{equation}\label{eq:lifted-H}
		\tilde H_j
		= \frac{1}{2}\sum_{i=1}^n g^{ii}_j(\boldsymbol{q})\, p_i^2
		+ V_j(\boldsymbol{q})\, p_u^2, \quad j=1,\ldots,n,\qquad \tilde{H}_{n+1}=p_u^2.
	\end{equation}
	
%%The integrals \eqref{eq:lifted-H} coincide with the Riemannian Eisenhart lifts associated with the Riemannian version of metric	\eqref{eq:3.3} (no $du dt$ term and vanishing vector potential $\boldsymbol{A}=0$).
	
In particular, $\tilde H_1$ is the geodesic Hamiltonian of the Riemannian Eisenhart lift of the natural system with metric $g_{ij}$ and potential $V_1$ (Riemannian wherever $V_1>0$), while the remaining $\tilde H_j$ are higher Stäckel integrals corresponding to Killing tensors of the same lifted metric.
\end{theorem}

\begin{proof}
	For Stäckel-separable systems, the contravariant metric and potentials admit
	the representations
	\begin{equation}\label{eq:metric-stackel}
		g^{kk}_j(\boldsymbol{q})=\bigl[(S^{(n)})^{-1}\bigr]_{j}^{k},
	\end{equation}
	and
	\begin{equation}\label{eq:potential-stackel}
		V_j(\boldsymbol{q})=\sum_{i=1}^n g^{ii}_j(\boldsymbol{q})\, W_i(q^i),
	\end{equation}
	see \cite{Benenti1997}.
	
	Since $S^{(n)}$ is invertible, the inverse of the lifting matrix
	\eqref{eq:lifting-matrix} is obtained via the Schur complement.
	For a block matrix
	\begin{equation}\label{eq:block-matrix}
		M=\begin{pmatrix} A & B \\ C & D \end{pmatrix},
	\end{equation}
	where $A$ is invertible and whose Schur complement
	\begin{equation}\label{eq:schur-complement}
		S_D = D - C A^{-1} B
	\end{equation}
	is also invertible, the inverse of $M$ is given by
	\begin{equation}\label{eq:schur-inverse}
		M^{-1}
		=
		\begin{pmatrix}
			A^{-1} + A^{-1} B S_D^{-1} C A^{-1}
			& -A^{-1} B S_D^{-1} \\
			- S_D^{-1} C A^{-1}
			& S_D^{-1}
		\end{pmatrix}.
	\end{equation}
	
	In the present case $A=S^{(n)}$, $B=-\boldsymbol{W}$,
	$C=\boldsymbol{0}^T$, and $D=1$, so that $S_D=1$. Substitution into
	\eqref{eq:schur-inverse} yields
	\begin{equation}\label{eq:lifted-inverse}
		[\mathcal{L}^{(n+1)}]^{-1}
		=
		\begin{pmatrix}
			(S^{(n)})^{-1} & (S^{(n)})^{-1}\boldsymbol{W} \\
			\boldsymbol{0}^T & 1
		\end{pmatrix}.
	\end{equation}
	
	The Stäckel lift defines a separable Hamiltonian system on $T^*\tilde Q$,
	$\dim\tilde Q=n+1$. The $j$-th lifted Hamiltonian is
	\begin{equation}\label{eq:H-tilde-def}
		\tilde H_j
		= \frac{1}{2}\sum_{i=1}^n [\mathcal{L}^{-1}]_{ji}\, p_i^2
		+ [\mathcal{L}^{-1}]_{j,n+1}\, p_u^2 .
	\end{equation}
	
	From \eqref{eq:lifted-inverse} and \eqref{eq:metric-stackel} we have
	\begin{equation}
		[\mathcal{L}^{-1}]_{ji}=g^{ii}_j(\boldsymbol{q}),
	\end{equation}
	while using \eqref{eq:potential-stackel} gives
	\begin{equation}
		[\mathcal{L}^{-1}]_{j,n+1}
		= \sum_{k=1}^n g^{kk}_j(\boldsymbol{q})\, W_k(q^k)
		= V_j(\boldsymbol{q}).
	\end{equation}
	Substituting into \eqref{eq:H-tilde-def} yields \eqref{eq:lifted-H}, completing
	the proof.
\end{proof}

\begin{thebibliography}{99}
\bibitem[Benent]{Benenti1997} S. Benenti. Intrinsic characterization of the variable separation in the Hamilton-Jacobi equation. \textit{J. Math. Phys.} \textbf{38}, 6578--6602 (1997).
\end{thebibliography}
\end{document}
